% Empirical paper (target: ICPC 2027 RENE, IEEE 10+2 pages). Venue not final.
\documentclass[10pt,conference]{IEEEtran}
\usepackage[T1]{fontenc}
\usepackage{booktabs,graphicx,xcolor,url,amsmath}
\usepackage[hidelinks]{hyperref}
\newcommand{\dROIntrude}{+9.1}
\newcommand{\dROIntrudeLo}{+5.0}
\newcommand{\dROIntrudeHi}{+13.1}
\newcommand{\dROCg}{+12.3}
\newcommand{\dROCgLo}{+6.8}
\newcommand{\dROCgHi}{+17.0}
\newcommand{\dRFIntrude}{+1.8}
\newcommand{\dRFIntrudeLo}{-0.9}
\newcommand{\dRFIntrudeHi}{+5.2}
\newcommand{\dRFCg}{+5.3}
\newcommand{\dRFCgLo}{+1.2}
\newcommand{\dRFCgHi}{+9.7}
\newcommand{\dMRRIntrude}{+5.9}
\newcommand{\dMRRIntrudeLo}{+2.7}
\newcommand{\dMRRIntrudeHi}{+9.3}
\newcommand{\dMRRCg}{+9.0}
\newcommand{\dMRRCgLo}{+4.8}
\newcommand{\dMRRCgHi}{+13.4}
\newcommand{\ROnejinx}{0.901}
\newcommand{\RFivejinx}{0.971}
\newcommand{\ROneintrude}{0.810}
\newcommand{\RFiveintrude}{0.953}
\newcommand{\ROnecg}{0.778}
\newcommand{\RFivecg}{0.918}
\newcommand{\ROneall}{0.906}
\newcommand{\RFiveall}{0.974}
\newcommand{\ROnehaiku}{0.898}
\newcommand{\RFivehaiku}{0.971}
\newcommand{\ROnelinks}{0.839}
\newcommand{\RFivelinks}{0.924}
\newcommand{\FARjinx}{2.5}
\newcommand{\FARintrude}{4.6}
\newcommand{\nStrategies}{66}
\newcommand{\nSigRFive}{0}
\newcommand{\nComparedRFive}{65}
\newcommand{\nPairs}{1867}
\newcommand{\nSelf}{1525}
\newcommand{\nCross}{342}
\newcommand{\nSameAgent}{228}
\newcommand{\nCrossAgent}{114}
\newcommand{\nNeg}{2965}
\newcommand{\nPRs}{76{,}209}

\newcommand{\todo}[1]{\textcolor{red}{[TODO: #1]}}
\begin{document}
\title{Duplicate Pull Requests from Coding Agents:\\A Live-Queue Evaluation of \nStrategies{} Detection Strategies}
\author{\IEEEauthorblockN{\todo{authors}}\IEEEauthorblockA{\todo{affiliations}}}
\maketitle

\begin{abstract}
Coding agents now open pull requests (PRs) at scale, and many of them duplicate work that is already open: in one hand-labelled sample, 23.7\% of failed agent PRs were duplicates~\cite{ehsani2026fail}. Duplicate-PR detectors were built and evaluated on human PRs, scored against random pairs of PRs. We ask how well they work on agent PRs when scored the way they are used: at the moment a new PR is opened, against the PRs actually open in that repository (a \emph{live queue}). On AIDev~\cite{li2026aidev} we build \nPairs{} live-queue duplicate pairs (\nCross{} between different authors) and \nNeg{} PRs with no duplicate, and compare \nStrategies{} strategies: lexical, code-overlap, link, embedding, existing duplicate-PR bots, learned combinations, cross-encoders, and LLM judges (Qwen2.5 1.5B--14B, Claude Haiku 4.5). A cheap learned combination of text, code, link, and timing signals (Jinx) ranks the true duplicate first for \ROnejinx{} of cross-author pairs versus \ROneintrude{} for the INTRUDE-style detector and \ROnecg{} for the best existing bot, with \FARjinx\% versus \FARintrude\% false alarms. No LLM judge or re-ranker improves on it, and none of the \nComparedRFive{} strategies beats the classic detector on Recall@5 with a confidence interval above zero. Random-pair scoring ranks strategies differently from live-queue scoring, and duplicates that share no files remain unsolved.
\end{abstract}

\section{Introduction}
\todo{Motivation: agent PR volume; duplicates in failure studies~\cite{ehsani2026fail,alam2026unmerged,nakashima2026rejected,abujadallah2026rejection}; maintainers' cost.}
\todo{Gap: no detector evaluated on agent PRs; random-pair negatives~\cite{ren2019intrude}; analogy to bug reports~\cite{zhang2023howfar}.}
\todo{Contributions: (1) live-queue benchmark on AIDev; (2) \nStrategies-strategy comparison incl.\ LLM judges; (3) duplicate taxonomy and per-class results; (4) evidence that random-pair scoring misleads; (5) open tool and artifact.}

\section{Research Questions}
\begin{itemize}
\item \textbf{RQ1} How well do existing duplicate-PR detectors and bots work on agent PRs under live-queue scoring?
\item \textbf{RQ2} Do embeddings, cross-encoders, or LLM judges beat cheap combined signals?
\item \textbf{RQ3} Which kinds of duplicates (taxonomy) remain hard?
\item \textbf{RQ4} Does random-pair scoring rank strategies the same way as live-queue scoring?
\end{itemize}

\section{Benchmark}
\todo{AIDev v5 (CC BY 4.0, cutoff 31 Mar 2026); silver labels from ``duplicate of / superseded by / in favor of \#N'' comments; live queue = PRs open in the repo when the new PR was created; \nPRs{} PRs; median queue 40, p90 314; \nSelf{} self-replacement, \nCross{} cross-author (\nSameAgent{} same agent, \nCrossAgent{} cross-agent); \nNeg{} no-duplicate queries.}
\todo{Taxonomy: A code twin, B partial overlap, C same files/different code, D no shared files; same issue/different code; same code/different description.}
\todo{Label audit result (precision of silver labels) -- pending.}

\section{Strategies}
\todo{Families and exact definitions; re-implemented bots and what was run as-is; training protocol: repository-grouped 5-fold CV.}

\section{Evaluation Protocol}
\todo{Recall@1/5, MRR; false-alarm rate at a threshold fixed on training folds; repository-cluster bootstrap (2000 resamples); paired differences; pre-registered kill rule.}

\section{Results}
\begin{table}[t]\centering\caption{Cross-author pairs ($n=\nCross$). False alarms on \nNeg{} PRs with no duplicate.}\label{tab:main}
\resizebox{\columnwidth}{!}{\begin{tabular}{lcccc}\toprule
Strategy & R@1 & R@5 & MRR & False alarms \% [95\% CI] \\\midrule
Jinx (ours) & 0.901 & 0.971 & 0.932 & 2.5 [1.9, 3.3] \\
All signals, LR & 0.906 & 0.974 & 0.935 & 1.6 [1.1, 2.1] \\
Jinx + BGE rerank + Qwen-7B & 0.901 & 0.974 & 0.931 & 2.8 [2.2, 3.5] \\
Claude Haiku 4.5 + Jinx & 0.898 & 0.971 & 0.932 & 1.8 [1.3, 2.3] \\
Qwen2.5-14B + Jinx & 0.871 & 0.971 & 0.917 & 2.9 [2.3, 3.7] \\
Claude Haiku 4.5 alone & 0.842 & 0.971 & 0.902 & 5.4 [4.6, 6.1] \\
INTRUDE + time~\cite{ren2019intrude,wang2019time} & 0.810 & 0.953 & 0.873 & 4.6 [3.7, 5.7] \\
contrib-guard (bot) & 0.778 & 0.918 & 0.843 & 11.7 [8.9, 14.6] \\
Links only & 0.839 & 0.924 & 0.882 & 3.8 [3.0, 4.8] \\
TF-IDF~\cite{li2017detecting} & 0.757 & 0.921 & 0.832 & 9.4 [7.8, 10.9] \\
BGE-small embedding & 0.763 & 0.915 & 0.830 & 6.4 [5.3, 7.5] \\
\bottomrule\end{tabular}
}\end{table}
\begin{table}[t]\centering\caption{Recall@1 per relation and duplicate class (cross-author unless noted).}\label{tab:classes}
\resizebox{\columnwidth}{!}{\begin{tabular}{lrccccc}\toprule
Group & $n$ & \cellcolor{jinxrow}Jinx & All LR & Haiku+Jinx & INTRUDE+t & contrib-guard \\\midrule
same agent & 207 & \cellcolor{jinxrow}0.923 & 0.928 & 0.928 & 0.913 & \textbf{0.932} \\
cross-agent & 106 & \cellcolor{jinxrow}0.934 & 0.925 & \textbf{0.953} & 0.915 & 0.877 \\
A: code twin & 132 & \cellcolor{jinxrow}\textbf{0.977} & 0.970 & \textbf{0.977} & 0.970 & 0.955 \\
B: partial overlap & 146 & \cellcolor{jinxrow}\textbf{0.932} & \textbf{0.932} & \textbf{0.932} & 0.918 & 0.904 \\
C: same files & 10 & \cellcolor{jinxrow}0.700 & 0.700 & \textbf{0.800} & 0.700 & 0.700 \\
D: no shared files & 25 & \cellcolor{jinxrow}0.720 & 0.760 & 0.800 & 0.680 & \textbf{0.840} \\
no link at all & 118 & \cellcolor{jinxrow}0.839 & 0.847 & \textbf{0.856} & 0.805 & \textbf{0.856} \\
self-replacement & 1454 & \cellcolor{jinxrow}0.891 & 0.891 & -- & \textbf{0.901} & 0.878 \\
\bottomrule\end{tabular}
}\end{table}
\textbf{RQ1.} Against the INTRUDE-style detector, Jinx improves cross-author Recall@1 by \dROIntrude{} points [\dROIntrudeLo, \dROIntrudeHi] and MRR by \dMRRIntrude{} [\dMRRIntrudeLo, \dMRRIntrudeHi], but Recall@5 only by \dRFIntrude{} [\dRFIntrudeLo, \dRFIntrudeHi]: the pre-registered Recall@5 criterion is not met. Against the best existing bot, Recall@5 improves by \dRFCg{} [\dRFCgLo, \dRFCgHi].
\textbf{RQ2.} \todo{LLM judges and rerankers: no CI-backed gain over Jinx; cost and latency.}
\textbf{RQ3.} \todo{Classes C and D; no-link pairs; same code/different description.}
\textbf{RQ4.} \todo{Kendall $\tau$ between random-pair AUC and live-queue Recall@1; strategies whose rank flips.}

\section{Threats to Validity}
\todo{Silver labels; link bias (links-only rule reaches \ROnelinks{} Recall@1); no agent--human pairs; re-implemented bots; single dataset.}

\section{Related Work}
Classic duplicate-PR detection~\cite{li2017detecting,yu2018dataset,ren2019intrude,wang2019time,jiang2023duphunter,liao2024dupllm,galbraith2025duplicate,benchekaya2026slr}; competing and redundant PRs~\cite{zhang2018competing,li2022redundancy}; agent-PR studies~\cite{ehsani2026fail,xu2026conflicts,sarkar2026before,khemissi2026humans}; duplicate bug reports~\cite{zhang2023howfar,cupid2023}. \todo{Commercial review tools and open bots (incl.\ doppelg\"anger, DupeHound).}

\section{Conclusion}
\todo{}

\section*{Data Availability}
\todo{Zenodo DOI for benchmark, code, and raw results.}
\bibliographystyle{IEEEtran}
\bibliography{../refs}
\end{document}
